\ifdefined\gkver\else\def\gkver{student}\fi
\documentclass[\gkver]{gaokaozhenti}
\begin{document}

\chapter{1988年上海}

\begin{enumerate}
\item
%% number: 1
%% paperName: 1988年上海高考第 1 题 4 分
%% typeId: 1
%% type: 单选题
%% chapter: 光学
%% point: 光电效应
%% method: 无
%% score: 4
%% degree: 850
%% duplicateId: 0
%% body:
关于光电效应的下述说法中，正确的是\xzanswer{C}
\fourchoices[answer=C]
{光电子的最大初动能随着入射光强度的增大而增大}
{只要入射光的强度足够大或照射的时间足够长，就一定能产生光电效应}
{任何一种金属都有一个极限频率，低于这个频率的光不能产生光电效应}
{在光电效应中，光电流的强度与入射光强度无关}
%% answer: C
%% memo:
\memoanswer{
本题原卷及材料中未提供详解，待补充。
}

\item
%% number: 2
%% paperName: 1988年上海高考第 2 题 4 分
%% typeId: 1
%% type: 单选题
%% chapter: 光学
%% point: 光的折射
%% method: 无
%% score: 4
%% degree: 650
%% duplicateId: 0
%% body:
两个平面镜互成直角，入射光线 AB 经过两次反射后的反射光线为 CD。今以两平面镜的交线为轴，将镜转动 $10^{\circ}$，两平面镜仍保持直角，在入射光 AB 保持不变的情况下，经过两次反射后，反射光线为 C$'$D$'$，则 C$'$D$'$ 与 CD\xzanswer{A}
\onepicture[w=4.2cm]{\includesvg{figs/02}}
\fourchoices[answer=A]
{不相交，同向平行}
{不相交，反向平行}
{相交成 $20^{\circ}$ 角}
{相交成 $40^{\circ}$ 角}
%% answer: A
%% memo:
\memoanswer{
本题原卷及材料中未提供详解，待补充。
}

\item
%% number: 3
%% paperName: 1988年上海高考第 3 题 4 分
%% typeId: 1
%% type: 单选题
%% chapter: 电磁感应
%% point: 楞次定律推广
%% method: 无
%% score: 4
%% degree: 650
%% duplicateId: 0
%% body:
位于载流长直导线近旁的两根平行铁轨 A 和 B，与长直导线平行且在同一水平面上，在铁轨 A、B 上套有两段可以自由滑动的导体 CD 和 EF，如图所示。若用力使导体 EF 向右运动，则导体 CD 将\xzanswer{B}
\onepicture[w=6.5cm]{\includesvg{figs/03}}
\fourchoices[answer=B]
{保持不动}
{向右运动}
{向左运动}
{先向右运动，后向左运动}
%% answer: B
%% memo:
\memoanswer{
本题原卷及材料中未提供详解，待补充。
}

\item
%% number: 4
%% paperName: 1988年上海高考第 4 题 4 分
%% typeId: 1
%% type: 单选题
%% chapter: 电磁感应
%% point: 右手定则
%% method: 无
%% score: 4
%% degree: 650
%% duplicateId: 0
%% body:
如图为一矩形线圈，绕与匀强磁场垂直的中心轴 OO$'$ 按顺时针方向旋转。引出线的两端各与互相绝缘的半圆铜环连接，两个半圆环分别与固定电刷 A、B 滑动接触，电刷间接有电阻 $R$，如图所示。在线圈转动过程中，通过 $R$ 的电流\xzanswer{C}
\onepicture[w=5.5cm]{\includesvg{figs/04}}
\fourchoices[answer=C]
{大小和方向都不断变化}
{大小和方向都不变}
{大小不断变化，方向从 A$\to$R$\to$B}
{大小不断变化，方向从 B$\to$R$\to$A}
%% answer: C
%% memo:
\memoanswer{
本题原卷及材料中未提供详解，待补充。
}

\item
%% number: 5
%% paperName: 1988年上海高考第 5 题 4 分
%% typeId: 1
%% type: 单选题
%% chapter: 曲线运动
%% point: 人造地球卫星
%% method: 无
%% score: 4
%% degree: 650
%% duplicateId: 0
%% body:
设人造地球卫星绕地球作匀速圆周运动，卫星离地面越高，则卫星的\xzanswer{D}
\fourchoices[answer=D]
{速度越大}
{角速度越大}
{向心加速度越大}
{周期越长}
%% answer: D
%% memo:
\memoanswer{
本题原卷及材料中未提供详解，待补充。
}

\item
%% number: 6
%% paperName: 1988年上海高考第 6 题 4 分
%% typeId: 1
%% type: 单选题
%% chapter: 功和能
%% point: 功率
%% method: 无
%% score: 4
%% degree: 650
%% duplicateId: 0
%% body:
设在平直公路上以一般速度行驶的自行车，所受阻力约为车、人总重量的 0.02 倍，则骑车人的功率最接近于\xzanswer{A}
\fourchoices[answer=A]
{$10^{-1}\UkW$}
{$10^{-3}\UkW$}
{$1\UkW$}
{$10\UkW$}
%% answer: A
%% memo:
\memoanswer{
本题原卷及材料中未提供详解，待补充。
}

\item
%% number: 7
%% paperName: 1988年上海高考第 7 题 4 分
%% typeId: 1
%% type: 单选题
%% chapter: 气体性质
%% point: 热力学第一定律
%% method: 无
%% score: 4
%% degree: 650
%% duplicateId: 0
%% body:
关于物体内能及其变化，下列说法中正确的是\xzanswer{B}
\fourchoices[answer=B]
{物体的温度改变时，其内能必定改变}
{物体对外作功，其内能不一定改变；向物体传递热量，其内能不一定改变}
{对物体作功，其内能必定改变；物体向外传出一定热量，其内能必定改变}
{若物体与外界不发生热交换，则物体的内能一定不改变}
%% answer: B
%% memo:
\memoanswer{
本题原卷及材料中未提供详解，待补充。
}

\item
%% number: 8
%% paperName: 1988年上海高考第 8 题 4 分
%% typeId: 1
%% type: 单选题
%% chapter: 气体性质
%% point: 玻意耳定律
%% method: 无
%% score: 4
%% degree: 650
%% duplicateId: 0
%% body:
右图是一端封闭，一端开口，截面均匀的 U 形玻璃管 AB，管内灌有水银。两水银面高度相等，闭管 A 内封有一定质量的气体，气体压强为 $72\UcmHg$。今将管 B 的开口端接到一抽气机，抽尽管 B 内水银面上的空气，结果两水银面产生 $18\Ucm$ 的高度差，问管 A 内原来气体柱的长度是\xzanswer{C}
\onepicture[w=2.6cm]{\includesvg{figs/08}}
\fourchoices[answer=C]
{$18\Ucm$}
{$12\Ucm$}
{$3\Ucm$}
{$6\Ucm$}
%% answer: C
%% memo:
\memoanswer{
本题原卷及材料中未提供详解，待补充。
}

\item
%% number: 9
%% paperName: 1988年上海高考第 9 题 5 分
%% typeId: 2
%% type: 多选题
%% chapter: 原子和原子核
%% point: 天然放射性
%% method: 无
%% score: 5
%% degree: 650
%% duplicateId: 0
%% body:
可以被电场加速或偏转的射线是\xzanswer{AC}
\fourchoices[answer=AC]
{阴极射线}
{伦琴射线}
{$\alpha$ 射线}
{$\gamma$ 射线}
%% answer: AC
%% memo:
\memoanswer{
本题原卷及材料中未提供详解，待补充。
}

\item
%% number: 10
%% paperName: 1988年上海高考第 10 题 5 分
%% typeId: 1
%% type: 单选题
%% chapter: 原子和原子核
%% point: 人工核反应
%% method: 无
%% score: 5
%% degree: 650
%% duplicateId: 0
%% body:
右图为查德威克实验示意图，由天然放射性元素钋（Po）放出的 $\alpha$ 射线轰击铍时会产生粒子流 A，用粒子流 A 轰击石蜡时，会打出粒子流 B，经研究知道\xzanswer{A}
\onepicture[w=9cm]{\includesvg{figs/10}}
\fourchoices[answer=A]
{A 为中子，B 为质子}
{A 为质子，B 为中子}
{A 为 $\gamma$ 射线，B 为中子}
{A 为中子，B 为 $\gamma$ 射线}
%% answer: A
%% memo:
\memoanswer{
本题原卷及材料中未提供详解，待补充。
}

\item
%% number: 11
%% paperName: 1988年上海高考第 11 题 5 分
%% typeId: 2
%% type: 多选题
%% chapter: 电场
%% point: 电场线
%% method: 无
%% score: 5
%% degree: 650
%% duplicateId: 0
%% body:
在电场中，任意取一条电场线，电场线上的 a、b 两点相距为 $d$，则\xzanswer{BD}
\onepicture[w=5cm]{\includesvg{figs/11}}
\fourchoices[answer=BD]
{a 点场强一定大于 b 点场强}
{a 点电势一定高于 b 点电势}
{a、b 两点间电势差一定等于 $Ed$（$E$ 为 a 点的场强）}
{a、b 两点间电势差等于单位正电荷由 a 点沿任意路径移到 b 点的过程中电场力做的功}
%% answer: BD
%% memo:
\memoanswer{
本题原卷及材料中未提供详解，待补充。
}

\item
%% number: 12
%% paperName: 1988年上海高考第 12 题 5 分
%% typeId: 2
%% type: 多选题
%% chapter: 运动和力
%% point: 动量守恒定律
%% method: 无
%% score: 5
%% degree: 650
%% duplicateId: 0
%% body:
质量为 $m$ 的小球 A，沿光滑水平面以 $v_{0}$ 的速度与质量为 $2m$ 的静止小球 B 发生正碰。碰撞后，A 球的动能变为原来的 $\frac{1}{9}$，那么小球 B 的速度可能是\xzanswer{AB}
\fourchoices[answer=AB]
{$\frac{1}{3}v_{0}$}
{$\frac{2}{3}v_{0}$}
{$\frac{4}{9}v_{0}$}
{$\frac{5}{9}v_{0}$}
%% answer: AB
%% memo:
\memoanswer{
本题原卷及材料中未提供详解，待补充。
}

\item
%% number: 13
%% paperName: 1988年上海高考第 13 题 5 分
%% typeId: 2
%% type: 多选题
%% chapter: 运动和力
%% point: 连接体
%% method: 整体与隔离
%% score: 5
%% degree: 650
%% duplicateId: 0
%% body:
两重迭在一起的滑块，置于固定的、倾角为 $\theta$ 的斜面上，如图所示，滑块 A、B 的质量分别为 $M$、$m$，A 与斜面间的滑动摩擦系数为 $\mu_{1}$，B 与 A 之间的滑动摩擦系数为 $\mu_{2}$。已知两滑块都从静止开始以相同的加速度从斜面滑下，滑块 B 受到的摩擦力\xzanswer{BC}
\onepicture[w=6.5cm]{\includesvg{figs/13}}
\fourchoices[answer=BC]
{等于零}
{方向沿斜面向上}
{大小等于 $\mu_{1}mg\cos\theta$}
{大小等于 $\mu_{2}mg\cos\theta$}
%% answer: BC
%% memo:
\memoanswer{
本题原卷及材料中未提供详解，待补充。
}

\item
%% number: 14
%% paperName: 1988年上海高考第 14 题 4 分
%% typeId: 3
%% type: 填空题
%% chapter: 气体性质
%% point: 气体图象
%% method: 无
%% score: 4
%% degree: 650
%% duplicateId: 0
%% body:
一定质量的理想气体，经历一膨胀过程，这过程可以用 $p-V$ 图上的直线 ABC 来表示。如图所示，在 A、B、C 三个状态上，气体的温度 $T_{A}$ \tkanswer{等于} $T_{C}$、$T_{B}$ \tkanswer{大于} $T_{A}$。（填“大于”、“等于”、“小于”）
\onepicture[align=right, w=5cm]{\includesvg{figs/14}}
%% answer: 等于，大于
\jdanswer{
等于，大于
}
%% memo:
\memoanswer{
本题原卷及材料中未提供详解，待补充。
}

\item
%% number: 15
%% paperName: 1988年上海高考第 15 题 4 分
%% typeId: 3
%% type: 填空题
%% chapter: 光学
%% point: 光的折射
%% method: 无
%% score: 4
%% degree: 650
%% duplicateId: 0
%% body:
如图所示，一束光线由某种媒质射入空气，当入射光线与界面的夹角为 $60^{\circ}$ 时，折射光线恰与反射光线垂直，该媒质的临界角等于 \tkanswer{$\sin^{-1}\frac{\sqrt{3}}{3}$}。（用反三角函数表示）
\onepicture[align=right, w=5cm]{\includesvg{figs/15}}
%% answer: \sin^{-1}\frac{\sqrt{3}}{3}
\jdanswer{
$\sin^{-1}\frac{\sqrt{3}}{3}$
}
%% memo:
\memoanswer{
本题原卷及材料中未提供详解，待补充。
}

\item
%% number: 16
%% paperName: 1988年上海高考第 16 题 4 分
%% typeId: 3
%% type: 填空题
%% chapter: 机械振动
%% point: 振动图象
%% method: 无
%% score: 4
%% degree: 850
%% duplicateId: 0
%% body:
甲乙两人先后观察同一弹簧振子在竖直方向上下振动的情况。
\begin{enumerate}
	\item
	甲开始观察时，振子正好在平衡位置并向下运动。试画出甲观察到的弹簧振子的振动图象。已知经过 $1\Us$ 后，振子第一次回到平衡位置。振子振幅为 $5\Ucm$（设平衡位置上方为正方向，时间轴上每格代表 $0.5\Us$）。
	\item
	乙在甲观察 $3.5\Us$ 后，开始观察并记时，试画出乙观察到的弹簧振子的振动图象。
\end{enumerate}
\drawpicanswer{\onepicture[align=right, w=6cm]{\includesvg{figs/16}}}{\onepicture[align=right, w=9cm]{\includesvg{figs/16_an}}}
%% answer: 见图
\jdanswer{
\begin{enumerate}
	\item
	见图

	\item
	见图

\end{enumerate}
}
%% memo:
\memoanswer{
本题原卷及材料中未提供详解，待补充。
}

\item
%% number: 17
%% paperName: 1988年上海高考第 17 题 4 分
%% typeId: 3
%% type: 填空题
%% chapter: 电路
%% point: 电功电功率
%% method: 无
%% score: 4
%% degree: 650
%% duplicateId: 0
%% body:
有位同学用以下方法测定一架收录机的功率：合上收录机的电键，把音量控制在适当范围，再把家中所有的其他用电器断开，然后观察自己家中使用的电度表（小火表）的转盘速度，发现转盘转过两圈所需时间为 $150\Us$。该电度表上标明每千瓦时（1 度电）转 1200 圈。由此可以计算出收录机的功率为 \tkanswer{$40\UW$}。
%% answer: 40
\jdanswer{
$40\UW$
}
%% memo:
\memoanswer{
本题原卷及材料中未提供详解，待补充。
}

\item
%% number: 18
%% paperName: 1988年上海高考第 18 题 4 分
%% typeId: 3
%% type: 填空题
%% chapter: 原子和原子核
%% point: 质能方程
%% method: 无
%% score: 4
%% degree: 650
%% duplicateId: 0
%% body:
\ce{^{232}_{92}U}（原子量为 232.0372）衰变为 \ce{^{238}_{90}Th}（原子量为 228.0287）时，释放出 $\alpha$ 粒子（\ce{^{4}_{2}He} 的原子量为 4.0026），则在衰变过程中释放出的能量为 \tkanswer{$8.8\times10^{-13}\UJ$}。（$1\Uu = 1.66\times10^{-27}\Ukg$）
%% answer: 8.8\times10^{-13}
\jdanswer{
$8.8\times10^{-13}\UJ$
}
%% memo:
\memoanswer{
本题原卷及材料中未提供详解，待补充。
}

\item
%% number: 19
%% paperName: 1988年上海高考第 19 题 4 分
%% typeId: 3
%% type: 填空题
%% chapter: 功和能
%% point: 机械能守恒定律
%% method: 等效替代
%% score: 4
%% degree: 450
%% duplicateId: 0
%% body:
两个底面积都是 $S$ 的圆桶，放在同一水平面上，桶内装水，水面高度分别为 $h_{1}$ 和 $h_{2}$，如图所示。已知水的密度为 $\rho$。现把连接两桶的阀门打开，最后两桶水面高度相等，则这过程中重力做的功等于 \tkanswer{$\frac{1}{4}\rho gS(h_{1}-h_{2})^{2}$}。
\onepicture[align=right, w=4.5cm]{\includesvg{figs/19}}
%% answer: \frac{1}{4}\rho gS(h_{1}-h_{2})^{2}
\jdanswer{
$\frac{1}{4}\rho gS(h_{1}-h_{2})^{2}$
}
%% memo:
\memoanswer{
本题原卷及材料中未提供详解，待补充。
}

\item
%% number: 20
%% paperName: 1988年上海高考第 20 题 6 分
%% typeId: 3
%% type: 填空题
%% chapter: 磁场
%% point: 带电粒子在复合场中的运动
%% method: 无
%% score: 6
%% degree: 550
%% duplicateId: 0
%% body:
MN 为水平放置的两块平行金属板，板间距离为 $d$，两板间电势差为 $U$。当带电量为 $q$，质量为 $m$ 的正离子流以速度 $v_{0}$ 沿水平方向从两板左端的中央 O 点处射入，因受电场力作用，离子作曲线运动。偏向 M 板（重力忽略不计）。今在两板间加一匀强磁场，使从中央 O 点处射入的正离子流在两板间作直线运动。则磁场的方向是 \tkanswer{垂直纸面指向纸外}，磁感应强度 $B =$ \tkanswer{$\frac{U}{v_{0}d}$}。如果把上述磁场的区域扩大到金属板右侧的空间，则从平板间射出的正离子将在这磁场中作圆周运动，完成半个圆周所需的时间为 \tkanswer{$\frac{\pi m}{qB}$}。
\onepicture[align=right, w=5.5cm]{\includesvg{figs/20}}
%% answer: 垂直纸面指向纸外；\frac{U}{v_{0}d}；\frac{\pi m}{qB}
\jdanswer{
垂直纸面指向纸外；$\frac{U}{v_{0}d}$；$\frac{\pi m}{qB}$
}
%% memo:
\memoanswer{
本题原卷及材料中未提供详解，待补充。
}

\item
%% number: 21
%% paperName: 1988年上海高考第 21 题 12 分
%% typeId: 4
%% type: 实验题
%% chapter: 电路
%% point: 故障分析
%% method: 无
%% score: 12
%% degree: 650
%% duplicateId: 0
%% body:
在做“测定电池的电动势和内电阻”的实验时，有位同学按下图电路进行连接，他共用 6 根导线，即 aa$'$、bb$'$、cc$'$、dd$'$、b$'$e 及 df，由于混进了一根内部断开的导线，所以当他按下开关 $K$ 后，发现两个电表的指针都不偏转。他用万用表电压档测量 ab$'$ 间电压时，读数约为 $1.5\UV$（已知电池电动势约为 $1.5\UV$）。为了确定哪一根导线的内部是断开的，可以再用万用表的电压档测量 \tkanswer{aa$'$} 两点间的电压，如果约为 $1.5\UV$，则一定是 \tkanswer{aa$'$} 导线断开；如果是 \tkanswer{$0\UV$}，则一定是 \tkanswer{bb$'$} 导线断开。

另一同学按上图把线路连接正确后，通过改变滑动变阻器接触点的位置，测出了几组 $U$、$I$ 的数据，（见图）是根据 1、2、3、4、5、6 组的数据在 $U-I$ 图上画出的点子和相应的 $U-I$ 图线。根据这一图线，可求出电池的电动势 $\varepsilon =$ \tkanswer{$1.43\UV$}，内电阻 $r =$ \tkanswer{$0.98\UO$}。如果他不利用这一图线，而是只利用任意两组 $U$、$I$ 数据，那么当他选择 \tkanswer{4}、\tkanswer{5} 两组数据时，求出的 $\varepsilon$、$r$ 值误差最大。
\twopicture[labelA=1988上海21a, labelB=1988上海21b, align=right, wA=5cm, wB=8cm]
{\includesvg{figs/21a}}
{\includesvg{figs/21b}}
%% answer: （1）aa$'$，aa$'$，0，bb$'$ 或 bb$'$，bb$'$，0，aa$'$；（2）1.43，0.98，4、5
\jdanswer{
\begin{enumerate}
	\item
	aa$'$，aa$'$，0，bb$'$；或 bb$'$，bb$'$，0，aa$'$。

	\item
	1.43，0.98，4、5。

\end{enumerate}
}
%% memo:
\memoanswer{
本题原卷及材料中未提供详解，待补充。
}

\item
%% number: 22
%% paperName: 1988年上海高考第 22 题 5 分
%% typeId: 4
%% type: 实验题
%% chapter: 运动和力
%% point: 动量守恒定律
%% method: 无
%% score: 5
%% degree: 650
%% duplicateId: 0
%% body:
在做“碰撞中的动量守恒”实验中，入射小球的质量应 \tkanswer{大于} 被碰小球的质量（填“大于”、“等于”或“小于”）。可以用小球飞出的水平距离代表小球碰撞前后的速度是因为 \tkanswer{小球都在同一高度作平抛运动，落地的时间相同，速度与水平距离成正比}。
%% answer: 大于；小球都在同一高度作平抛运动，落地的时间相同，速度与水平距离成正比
\jdanswer{
\begin{enumerate}
	\item
	大于。

	\item
	小球都在同一高度作平抛运动，落地的时间相同，速度与水平距离成正比。

\end{enumerate}
}
%% memo:
\memoanswer{
本题原卷及材料中未提供详解，待补充。
}

\item
%% number: 23
%% paperName: 1988年上海高考第 23 题 5 分
%% typeId: 4
%% type: 实验题
%% chapter: 机械振动
%% point: 单摆
%% method: 无
%% score: 5
%% degree: 650
%% duplicateId: 0
%% body:
做“用单摆测定重力加速度”的实验，下述说法中正确的是\xzanswer{AB}
\fourchoices[answer=AB]
{如果有两个大小相同的铁球和木球（都有小孔）可供选择，则选用铁球作为摆球较好}
{单摆的偏角不要超过 $5^{\circ}$}
{为了便于改变摆线的长度，可将摆线的一头绕在铁架上端的圆杆上以代替铁夹}
{测量摆长时，应该用力拉紧摆线}
%% answer: AB
\jdanswer{
AB
}
%% memo:
\memoanswer{
本题原卷及材料中未提供详解，待补充。
}

\item
%% number: 24
%% paperName: 1988年上海高考第 24 题 12 分
%% typeId: 6
%% type: 计算题
%% chapter: 电路
%% point: 闭合电路欧姆定律
%% method: 极值
%% score: 12
%% degree: 650
%% duplicateId: 0
%% body:
如图所示电路：已知电源电动势 $E = 6.3\UV$，内电阻 $r = 0.5\UO$，固定电阻 $R_{1} = 2\UO$、$R_{2} = 3\UO$。$R_{3}$ 是阻值为 $5\UO$ 的滑动变阻器。按下电键 K，调节滑动变阻器的触点，求通过电源的电流范围。
\onepicture[align=right, w=5cm]{\includesvg{figs/24}}
%% answer: 2.1 A 到 3 A 之间
\jdanswer{
$2.1\UA$ 到 $3\UA$ 之间
}
%% memo:
\memoanswer{
\onepicture[align=right, w=4cm]{\includesvg{figs/24_an}}

设变阻器的触点把电阻 $R_{3}$ 分成 $R_{0}$ 和 $R_{3}-R_{0}$ 两部分，等效电路如图所示，电路的外电阻

$R=\frac{1}{\frac{1}{R_{2}+R_{0}}+\frac{1}{R_{1}+R_{3}-R_{0}}}=\frac{(5-R_{0}+2)(3+R_{0})}{10}=\frac{25-(R_{0}-2)^{2}}{10}$ ①

当 $R_{0}-2=0$ 即 $R_{0}=2\UO$ ②

外电阻最大，其大小为

$R_{\text{最大}}=\frac{(5-2+2)(3+2)}{10}\UO=2.5\UO$ ③

当外电阻最大时，通过电池的电流最小

$I_{\text{最小}}=\frac{E}{R_{\text{最大}}+r}=\frac{6.3}{2.5+0.5}\UA=2.1\UA$ ④

由①式，当 $R_{0}=R_{2}=5\UO$ ⑤

时，外电阻最小，其大小为

$R_{\text{最小}}=\frac{(5-5+2)(3+5)}{10}\UO=1.6\UO$ ⑥

当外电阻最小时，通过电池的电流最大

$I_{\text{最大}}=\frac{E}{R_{\text{最小}}+r}=\frac{6.3}{1.6+0.5}\UA=3\UA$ ⑦

电流的变化范围在 $2.1\UA$ 到 $3\UA$ 之间。
}

\item
%% number: 25
%% paperName: 1988年上海高考第 25 题 14 分
%% typeId: 6
%% type: 计算题
%% chapter: 磁场
%% point: 磁力矩
%% method: 无
%% score: 14
%% degree: 450
%% duplicateId: 0
%% body:
一具有固定转轴的矩形导线框 abcd，处在直线电流的磁场中，转轴与直导线平行，相距 $4r_{0}$，线框的 ab 和 cd 两边与转轴平行，长度都为 $5r_{0}$，bc 和 da 两边与转轴垂直，长度都为 $6r_{0}$，转轴通过这两条边的中点，如图所示。直导线中的电流向上。当导线框垂直于由直线电流与转轴构成的平面时，ab 边和 cd 边所在处的磁感应强度的大小都是已知值 $B_{0}$。
\begin{enumerate}
	\item
	若导线框以恒定的角速度 $\omega$ 绕固定轴转动，求当导线框转到上述位置时，线框中的感应电动势。
	\item
	若导线框转到上述位置时，框内电流强度为 $I$，方向为 abcda，求导线框处在这位置时，磁场对 ab 边和 cd 边的作用力以及这两个力对转轴的力矩。
\end{enumerate}
\onepicture[align=right, w=5cm]{\includesvg{figs/25}}
%% answer: （1）E=24B_{0}r_{0}^{2}\omega；（2）F=5Ir_{0}B_{0}，M=24B_{0}r_{0}^{2}I
\jdanswer{
\begin{enumerate}
	\item
	$E=24B_{0}r_{0}^{2}\omega$

	\item
	$F=5Ir_{0}B_{0}$，$M=24B_{0}r_{0}^{2}I$

\end{enumerate}
}
%% memo:
\memoanswer{
\onepicture[align=right, w=5cm]{\includesvg{figs/25_an1}}

（1）导线框绕固定轴转动时，ab、cd 两边都作切割磁感线运动。设导线框绕转轴按逆时针方向旋转，当导线框转到给定的位置时，ab 边和 cd 边速度的大小相等，这速度为

$v=3r_{0}\omega$ ①

速度的方向、ab 和 cd 边所在处磁场的方向如图所示。由于 ab 边与所在处的磁场方向垂直，速度方向与磁感线的夹角为 $\varphi$，故 ab 边中的感应电动势为

$E=B_{0}\cdot ab\cdot v\sin\varphi$ ②

而 $\sin\varphi=\cos\theta=\frac{4r_{0}}{\sqrt{(4r_{0})^{2}+(3r_{0})^{2}}}=0.8$ ③

已知，$ab=5r_{0}$，由①②式，得

$E_{ab}=B_{0}\times5r_{0}\times3r_{0}\omega\times0.8=12B_{0}r_{0}^{2}\omega$ ④

cd 边中的感生电动势与 ab 边中的感生电动势的大小相同，而且又是串联在一起的，所以这一瞬间导线框中的感生电动势是

$E=2E_{ab}=24B_{0}r_{0}^{2}\omega$ ⑤

\onepicture[align=right, w=5cm]{\includesvg{figs/25_an2}}

（2）当导线框位于给定位置时，导线 ab 所在处的磁感应强度的大小为 $B_{0}$，磁场的方向与通电导线 ab 垂直，如图所示。所以磁场对导线 ab 的作用力的大小为

$f_{ab}=B_{0}I\cdot ab=5Ir_{0}B_{0}$ ⑥

磁场对 cd 边的作用力的大小等于

$f_{cd}=5Ir_{0}B_{0}$ ⑦

$f_{ab}$ 对转轴的力矩

$M_{ab}=f_{ab}L_{2}=f_{ab}\cdot O'b\cdot\cos\theta=5Ir_{0}B_{0}\times3r_{0}\times0.8=12B_{0}r_{0}^{2}I$ ⑧

$f_{cd}$ 对转轴的力矩 $M_{cd}$ 的大小等于 $M_{ab}$，两个力矩都有使导线框作顺时针转动的趋势，故合力矩

$M=M_{ab}+M_{cd}=24B_{0}r_{0}^{2}I$ ⑨
}

\item
%% number: 26
%% paperName: 1988年上海高考第 26 题 15 分
%% typeId: 6
%% type: 计算题
%% chapter: 运动和力
%% point: 动量和能量
%% method: 无
%% score: 15
%% degree: 350
%% duplicateId: 0
%% body:
在光滑水平面上，有一质量 $m_{1}=20\Ukg$ 的小车，通过一根几乎不可伸长的轻绳与另一个质量为 $m_{2}=25\Ukg$ 的拖车相连接。一质量 $m_{3}=15\Ukg$ 的物体放在拖车的平板上。物体与平板间的滑动摩擦系数为 $\mu=0.20$。开始时，拖车静止，绳未拉紧（如图所示），小车以 $v_{0}=3\Ums$ 的速度向前运动。求：
\begin{enumerate}
	\item
	当 $m_{1}$、$m_{2}$、$m_{3}$ 以同一速度前进时，速度的大小；
	\item
	物体在拖车平板上移动的距离。（$g$ 取 $10\Umsq$）
\end{enumerate}
\onepicture[align=right, w=9cm]{\includesvg{figs/26}}
%% answer: （1）v=1\Ums；（2）\Delta s=0.33\Um
\jdanswer{
\begin{enumerate}
	\item
	$v=1\Ums$

	\item
	$\Delta s=0.33\Um$

\end{enumerate}
}
%% memo:
\memoanswer{
（1）在从绳开始拉紧到 $m_{1}$，$m_{2}$ 和 $m_{3}$ 以同一速度运动的这一过程中，$m_{1}$、$m_{2}$ 和 $m_{3}$ 这三个物体不受外力作用，它们的总动量保持不变。由动量守恒定律

$m_{1}v_{0}=(m_{1}+m_{2}+m_{3})v$ ①

式中 $v_{0}$ 为 $m_{1}$ 的初速度，$v$ 为三者一起运动的速度。由①式得

$v=\frac{m_{1}v_{0}}{m_{1}+m_{2}+m_{3}}=\frac{20\times3}{20+25+15}\Ums=1\Ums$ ②

（2）在绳拉紧的极短时间内，$m_{1}$ 和 $m_{2}$ 的相互作用可看作一种碰撞过程，在这过程中，$m_{3}$ 作用于 $m_{2}$ 的摩擦力可忽略不计。由动量守恒定律

$m_{1}v_{0}=(m_{1}+m_{2})v_{2}$ ③

$v_{2}=\frac{m_{1}v_{0}}{m_{1}+m_{2}}=\frac{20\times3}{20+25}\Ums=\frac{4}{3}\Ums$ ④

当 $m_{1}$ 和 $m_{2}$ 以速度 $v_{2}$ 运动时，要受到 $m_{3}$ 的摩擦力的作用。因克服摩擦力作功，$m_{1}$ 和 $m_{2}$ 的速度由 $v_{2}$ 减少到 $v$。由动能定理

$fs_{2}=\frac{1}{2}(m_{1}+m_{2})v_{2}^{2}-\frac{1}{2}(m_{1}+m_{2})v^{2}$ ⑤

式中 $s_{2}$ 是 $m_{1}$ 和 $m_{2}$ 一起向前移动的距离，而摩擦力 $f$ 为

$f=\mu m_{3}g$ ⑥

由⑤、⑥以及④式，可得

$s_{2}=\frac{(m_{1}+m_{2})(v_{2}^{2}-v^{2})}{2\mu m_{3}g}=\frac{(20+25)(\frac{16}{9}-1)}{2\times0.2\times15\times10}\Um=\frac{7}{12}\Um$ ⑦

作用于 $m_{3}$ 的摩擦力对 $m_{3}$ 作功，使其速度增加至 $v$。由动能定理

$fs_{3}=\frac{1}{2}m_{3}v^{2}$ ⑧

由⑥、⑧以及②式，可得

$s_{3}=\frac{v^{2}}{2\mu g}=\frac{1}{2\times0.2\times10}\Um=0.25\Um$ ⑨

$\Delta s=s_{2}-s_{3}=0.33\Um$ ⑩
}

\end{enumerate}

\end{document}
